\documentclass[10pt,a4paper]{amsart}
\usepackage[margin=19mm]{geometry}
\usepackage{amsmath,amssymb,mathtools}
\usepackage[scr=rsfs,cal=euler]{mathalfa}
\usepackage{xfrac}
\usepackage[unicode,hidelinks,bookmarks=false]{hyperref}
\newcommand{\ie}{i.e.\ }
\newcommand{\cf}{cf.\ }
\newcommand{\resp}{respectively\ }
\newcommand{\Subsection}{\subsection}
\providecommand{\nobreakdash}{\nobreak\hbox{-}}
\setlength{\emergencystretch}{2em}
\multlinegap0pt
\usepackage{mathtools}
\usepackage{amssymb}
\newtheorem{lemma}{Lemma}
  \newcommand{\C}{{\mathcal{C}}}
  \newcommand{\E}{{\mathcal{E}}}
  \newcommand{\J}{{\mathcal{J}}}
  \newcommand{\M}{{\mathcal{M}}}
  \newcommand{\N}{{\mathcal{N}}}
  \newcommand{\T}{{\mathcal{T}}}
\def\al{\alpha}
\newcommand{\bo}[1]{\mathbf{#1}} %bold math symbols
\def\de{\delta}
\newcommand{\fH}{{\mathfrak{H}}}
\newcommand{\fg}{{\mathfrak{g}}}
\newcommand{\fh}{{\mathfrak{h}}}
\newcommand{\foral}{\text{ for all }}
\def\la{\lambda}
\title{The reduced Hao--Ng isomorphism problem for non-degenerate product systems}
\author{Ioannis Apollon Paraskevas}
\date{Source: arXiv:2607.17824v1 (CC BY 4.0)}
\begin{document}
\maketitle

\noindent\emph{Source and licence.} The reduced Hao--Ng isomorphism problem for non-degenerate product systems. Version arXiv:2607.17824v1 (CC BY 4.0), distributed by arXiv under the Creative Commons Attribution 4.0 International licence, http://creativecommons.org/licenses/by/4.0/, as recorded in the arXiv licence metadata for this exact version and shown on the abstract page https://arxiv.org/abs/2607.17824v1. Full licence notice: \texttt{licenses/arxiv-2607.17824-CC-BY-4.0.txt}.\par\smallskip

\noindent\textit{Editorial scope.} Counted unit: the article's lemma L:invar (source0.tex lines 1369--1373). The article's complete original proof of it is delivered with it (source0.tex lines 1375--1410, one \texttt{proof} environment of 36 lines). No mathematics is written, repaired or re-derived: the statement and the proof are the article's own lines, in the article's own notation, and no retained line is substituted or re-flowed.  Retained uncounted as the source-local context the proof needs: lines 1012--1014; lines 1016--1020; lines 1275--1277; lines 1282--1307.  Nothing was omitted from any retained span, so the package carries no omission record.  No retained line contains a citation, so the wrapper carries no bibliography and no bibliography entry.  No retained line contains a figure environment, an image-inclusion command or drawing code, so no image is delivered and the package declares no asset dependency.  Editorial formatting only, in addition: an AMS \texttt{amsart} preamble that restates the article's own declarations and macros used by the retained lines, namely the environments \texttt{lemma} on separate counters, together with the standard packages those lines need.  The retained environments print in this document as Lemma 1 in order of appearance, whereas the article prints the same items with the numbering of its own full document.  The complete native source of the article as distributed in the arXiv e-print for this exact version is bundled unchanged as \texttt{sources/205598/source0.tex}.
\begin{equation}\label{Eq:Buss}
C_c(\fH,\C)\subseteq\N_{\de^{\rtimes}},
\end{equation}

\begin{equation}\label{Eq:avercomp}
\E_{\de^{\rtimes}}(f^* \ast h)= i_\C^r\left((f^*\ast h)(e_{\fH})\right)= i_\C^r\left(
\int_{\fH}
\al_{\fg^{-1}}\big(f(\fg)^*h(\fg)\big) d\mu(\fg)\right).
\end{equation}

\begin{equation}\label{Eq:appunit}
u_i(\fh)=f_i(\fh)\la_{e_P}(e_i) \foral \fh\in \fH,
\end{equation}

\begin{lemma}\label{L:local1}
With the aforementioned notation, for each fixed $u_i$ defined in \eqref{Eq:appunit} we have:
\begin{enumerate}
\item $u_i \in \N_{\de^\rtimes}$.
\item If $z\in\T_\la(X)\rtimes_{\al,\la}\fH$, then $z u_i \in \N_{\de^\rtimes}$.

\item If $z,w\in \T_\la(X)\rtimes_{\al,\la}\fH$, then
$(z u_i)^*(w u_i)\in\M_{\de^{\rtimes}}$ and
$\E_{\de^{\rtimes}}((z u_i)^*(w u_i))\in j(\T_\la(X))$.
In particular, the map
\[
\Theta_i \colon
(\T_\la(X)\rtimes_{\al,\la}\fH)\times
(\T_\la(X)\rtimes_{\al,\la}\fH)\to j(\T_\la(X))
\text{ such that }
\Theta_i(z,w)
:=
\E_{\de^{\rtimes}}((zu_i)^*(w u_i)),
\]
is sesquilinear and
$\|\Theta_i(z,w)\|
\leq
\|z\|\cdot \|w\|\cdot \|\E_{\de^\rtimes}(u_i^*u_i)\|$.

\end{enumerate}
\end{lemma}

\begin{lemma}\label{L:invar}
Let $\bo{x}, \bo{y} \in \J$ be constructible ideals of $P$. With the aforementioned notation, for each fixed $u_i$ defined in \eqref{Eq:appunit} we have that if $z\in \bo{K}_{\bo{x},\la_\ast}\rtimes_{\al,\la}\fH$ and
$w\in \bo{K}_{\bo{y},\la_\ast}\rtimes_{\al,\la}\fH$, then
$\E_{\de^{\rtimes}}((z u_i)^*(w u_i))\in j(\bo{K}_{\bo{x}\cap \bo{y},\la_\ast})$.
\end{lemma}

\begin{proof}
First suppose that $z\in C_c(\fH,\bo{K}_{\bo{x},\la_\ast})$ and
$w\in C_c(\fH,\bo{K}_{\bo{y},\la_\ast})$. We have
\[
(z \ast u_i)(\fg)=\int_{\fH}z(\fh) f_i(\fh^{-1}\fg) \al_{\fh}(\la_{e_P}(e_i))\,d\mu(\fh)
\]
and
\[
(w \ast u_i)(\fg)=\int_{\fH}w(\fh)f_i(\fh^{-1}\fg) \al_{\fh}(\la_{e_P}(e_i))\,d\mu(\fh).
\]
Since $\al_{\fh}(\la_{e_P}(A))=\la_{e_P}(A)$ for every $\fh\in \fH$ and $\bo{K}_{\bo{x},\la_\ast} \cdot \la_{e_P}(A) \subseteq \bo{K}_{\bo{x},\la_\ast}$, and similarly for $\bo{K}_{\bo{y},\la_\ast}$, we get $z \ast u_i\in C_c(\fH,\bo{K}_{\bo{x},\la_\ast})$ and $w \ast u_i\in C_c(\fH,\bo{K}_{\bo{y},\la_\ast})$.
By \eqref{Eq:Buss} and \eqref{Eq:avercomp} we get
\[
\E_{\de^{\rtimes}}((z \ast u_i)^* \ast (w\ast u_i))
=
j\left(\int_{\fH}\al_{\fh^{-1}}
\big((z \ast u_i)(\fh)^*(w \ast u_i)(\fh)\big)\,d\mu(\fh)\right).
\]
Since the closed subspace $[\bo{K}_{\bo{x},\la_\ast} \cdot \bo{K}_{\bo{y},\la_\ast} ]$ is $\al$-invariant the integral belongs to
\[
[\bo{K}_{\bo{x},\la_\ast} \cdot \bo{K}_{\bo{y},\la_\ast} ] \subseteq \bo{K}_{\bo{x}\cap \bo{y},\la_\ast}.
\]
For arbitrary
$z\in \bo{K}_{\bo{x},\la_\ast} \rtimes_{\al,\la}\fH$ and
$w\in \bo{K}_{\bo{y},\la_\ast} \rtimes_{\al,\la}\fH$ choose nets
$z_\nu\in C_c(\fH,\bo{K}_{\bo{x},\la_\ast})$ and $w_{\nu'}\in C_c(\fH,\bo{K}_{\bo{y},\la_\ast})$ with $z_\nu \xrightarrow{\nu}z$ and $w_{\nu'} \xrightarrow{\nu'} w$ in norm. By Lemma~\ref{L:local1} (iii) we have
\[
\|\Theta_i(z_\nu,w_{\nu'})-\Theta_i(z,w)\|
\xrightarrow{\nu,\nu'} 0.
\]
In particular, since $\Theta_i(z_\nu,w_{\nu'})\in j(\bo{K}_{\bo{x}\cap \bo{y},\la_\ast})$ we conclude that
\[
\Theta_i(z,w)=\E_{\de^{\rtimes}}((z u_i)^*(w u_i))\in j(\bo{K}_{\bo{x}\cap \bo{y},\la_\ast}),
\]
as required.
\end{proof}

\end{document}
